\documentclass{homework}
% \usepackage{graphicx} % Required for inserting images
\usepackage{amsmath}
\usepackage{gensymb}
\usepackage{subfig}
\usepackage{float}

\class{ECEN 5478}
\title{Midterm}
\author{Sean Jennings}
\date{October 31, 2023}

\begin{document}
\maketitle

\section{Problem 1}

\begin{letterenum}
    \item A twice continuously differentiable function $f: \R^n \to \R$ is
        $\mu$-strongly convex if and only if
        \begin{equation*}
            y^T \nabla^2 f(x) y \geq \mu \norm{y}^2
        \end{equation*}
        for all $x,y \in \R^n$ [Sidford Notes 2020]. If $y=0$ the inequality
        is trivially satisfied. Assume $y\neq 0$ and divide both sides by $\norm{y}^2$
        to obtain

        \begin{equation*}
            \frac{y^T \nabla^2 f(x) y}{\norm{y}^2} \geq \mu
        \end{equation*}
        The left hand side is the Rayleigh coefficient,
        and which is lower-bounded as follows [Beck14, Lemma 1.11]:
        \begin{equation*}
            \frac{y^T \nabla^2 f(x) y}{\norm{y}^2} \geq \lambda_{min} (\nabla^2 f(x)).
        \end{equation*}
        The preceding equation becomes an equality when $y$ is the eigenvector
        corresponding to $\lambda_{min}(\nabla^2 f(x))$ since
        \begin{equation*}
            y^T \nabla^2 f(x) y = y^T \lambda_{min}(\nabla^2 f(x)) y
                = \lambda_{min}(\nabla^2 f(x)) \norm{y}^2.
        \end{equation*}

        Since the lower bound of the Rayleigh coefficient is attained,
        $f$ is $\mu$-strongly convex if and only if $\lambda_{min} (\nabla^2f(x)) > 0$.

        For the least square's problem, $f_t$ is indeed twice continuously
        differentiable with
        \begin{equation*}
            \nabla f_t(x) = A^T(Ax - b_t)
        \end{equation*}
        and
        \begin{equation*}
            \nabla^2 f_t(x) = A^TA
        \end{equation*}
        Therefore, if $\lambda_{min}(A^TA) > 0$, we have $f_t$ is $\mu$-strongly
        convex with coefficient $\mu = \lambda_{min}(A^TA)$. Otherwise, $f_t$
        is not $\mu$-strongly convex.

    % \item We will derive the second-order characterization for
    %     strongly convex functions. The gradient inequality for a strongly
    %     convex function $f: \R^n \to \R$ is given by:
    %     \begin{equation*}
    %         f(y) \geq f(x) + \nabla f(x)^T (y-x) + \frac{\mu}{2} \norm{y - x}^2
    %     \end{equation*}
    %     for any $x,y\in\R^n$. Swapping $x$ and $y$, we also have
    %     \begin{equation*}
    %         f(x) \geq f(y) - \nabla f(y)^T (y-x) + \frac{\mu}{2} \norm{y - x}^2.
    %     \end{equation*}
    %     Adding these two equations yields the fact that the gradient of a $\mu$-strongly
    %     convex function is strictly monotone.
    %     \begin{equation*}
    %         (\nabla f(x) - \nabla f(y))^T (x - y) \geq \mu \norm{x - y}^2
    %     \end{equation*}
    %     Taking the change of variables $x = y + h$ for some $h\in\R^n$ and assuming
    %     that $f$ is twice continuously differentiable, we also have:
    %     \begin{align*}
    %         h^T (\nabla f(y + h) - \nabla f(y))
    %             &= h^T (\nabla f(y + h) - \nabla f(y) - \nabla^2 f(x) h) + h^T \nabla^2 f(y) h \\
    %             &\geq \norm{h} \norm{\nabla f(y + h) - \nabla f(y) - \nabla^2 f(y) h} + \lambda_{min} (\nabla^2 f(y)) \norm{h}^2,
    %     \end{align*}
    %     where the final inequality follows from the fact that the Rayleight quotient is bounded by the minimum and maximum
    %     eigenvalues [Lemma 1.1 Beck14].

    %     Noting that by the definition of differentiability, we have
    %     \begin{equation*}
    %         \lim_{\norm{h} \to 0} \frac{\norm{\nabla f(y + h) - \nabla f(y) - \nabla^2 f(y) h}}{\norm{h}} = 0,
    %     \end{equation*}

    %     dividing both sides by $\norm{h}^2$ and taking the limit as $\norm{h} \to 0$, 
    %     it follows that
    %     \begin{equation}
    %         \lambda_{min}(\nabla^2 f(x)) \geq 
    %     \end{equation}
    %     for all $x$.

    %     For the least square's problem at hand, we have:
    %     \begin{align*}
    %         \nabla (f_t(x)) = A^T (Ax - b_t) \\
    %         \nabla^2 (f_t(x)) = A^T A.
    %     \end{align*}
    %     Therefore, $f_t$ is $\mu$-strongly convex iff $\lambda_{min}(A^TA) \geq \mu$. 

    \item Since $\nabla f_t(x) = A^T(Ax - b_t)$, we have

    \begin{align*}
        \norm{\nabla f_t (x) - \nabla f_t(y)} &= 
           \norm{A^T(Ax - b_t) - A^T(Ay - b_t)} \\
           &= \norm{A^TA (x - y)} \\
           &\leq \norm{A^TA} \norm{x - y}
    \end{align*}
    Therefore, the gradient is Lipschitz with constant
    $L=\norm{A^TA}=|\lambda_{max}(A^TA)|$.

    % \item For a twice continuously differentiable function, $L$-smoothness is
    %     equivalent to the 2-norm of the hessian being bounded by $L$:
    %     \begin{equation*}
    %         \norm{\nabla^2 f(x)} \leq L
    %     \end{equation*}
    %     for all $x\in\R^n$ [Theorem 5.12, Beck17]. By definition of an 
    %     $L$-smooth function, we also have:
    %     \begin{equation*}
    %         \norm{\nabla f(x) - \nabla f(y)} \leq L \norm{x - y}
    %     \end{equation*}
    %     which is equivalent to the gradient being Lipschitz with constant $L$.

    %     Therefore, the gradient of $f_t$ is $L$-Lipschitz if and only if
    %     \begin{equation*}
    %         \norm{\nabla^2 f_t(x)} = \norm{A^TA} = |\lambda_{max}(A^TA)|
    %     \end{equation*}

    %     Since $A^TA$ is symmetric, it's norm is given by
    %     $\norm{A^TA} = \lambda_{max} (A^TA)$.

    \item Since the the optimization problem is unconstrained (the domain is $R^n$),
        we have
        \begin{equation}
            \nabla f_t(x_t^*) = A^T (A x_t^* - b_t) = 0
            \label{eq:stationarity}
        \end{equation}
        since stationarity is a necessary condition [Theorem 7.9 Beck14].
    
        Since we now assume that $f_t$ is $\mu$-strongly convex, we have
        $\lambda_{min} (A^TA) > 0$. Therefore, $A^TA$ is full-rank and
        invertible (note if we were to have $d < n$, we would have a contradiction,
        since $A^TA$ can have rank at most $d$). We can solve then solve the
        equation (\ref{eq:stationarity}) for $x_t^*$:
        \begin{equation*}
            x_t^* = (A^TA)^{-1} A^Tb_t
        \end{equation*}

        Then, the drift in the optimal solution as a function of time is related
        to the change in the $b_t$ by:
        \begin{align*}
            \sigma_t &= \norm{x_t^* - x_{t-1}^*} \\
                &= \norm{(A^TA)^{-1}A^T b_t - (A^TA)^{-1} A^T b_{t-1}} \\
                &\leq \norm{(A^TA)^{-1} A^T} \norm{b_t - b_{t-1}} \\
                &= \norm{(A^TA)^{-1} A^T} \xi_t
        \end{align*}

    \item We will use the online gradient descent algorithm defined as follows:
    \begin{equation}
        x_{t+1} = x_t + \alpha \nabla f_t(x_t)
    \end{equation}
    The tracking bound of the online gradient descent algorithm is given
    by:
    \begin{equation*}
        \norm{x_t - x_t^*} \leq \rho^t \norm{x_0 - x_0^*}
            + \sum_{i=0}^{t-1} \rho^{t - 1 - i} \sigma_i \\
    \end{equation*}
    where $\rho = \max\{|1 - \alpha L|, |1 - \alpha \mu\}$. This result follows
    from the tracking bound given in lecture [Module 2, Slide 13], by noting that
    the convexity, smoothness coefficients, and step size do not depend on time,
    allowing us to convert the product of distinct Lipschitz coefficient to the
    exponents $\rho^{t-1}$ and $\rho^{t-1-i}$.

    For the algorithm to converge and track the optimal solution, we must
    have $\rho < 1$. Therefore the step size should be chosen so that
    \begin{align*}
        | 1 - \alpha L | &< 1 \\
        -(1 - \alpha L) &< 1 \\
        1 - \alpha L &> -1 \\
        \alpha &< 2 / L
    \end{align*}
    and similarly, so that we also have that $\alpha < 2 / \mu$.
    But since $\mu < L$, we
    also have $2 / \mu > 2 / L$, and the upper bound on step-size to obtain
    convergence is determined solely by the smoothness coefficient:
    \begin{equation*}
        \alpha < \frac{2}{L}.
    \end{equation*}

    \item Given that $A_t$ is diagonal (and thus symmetric), we now have:
    \begin{align*}
        (A_t^TA_t)^{-1} &= \text{diag}([a_{0,t}^2, \dots, a_{n,t}^2]) \\
            &= \text{diag}([1/a_{0,t}^2, \dots, a_{n,t}^2]) \\
    \end{align*}
    and
    \begin{align*}
        (A_t^TA_t)^{-1} A^T &= \text{diag}([1/a_{0,t}, \dots, a_{n,t}]) \\
            &= A_t^{-1}
    \end{align*}
    Thus, our convergence bound from the previous problem becomes:
    \begin{align*}
        \norm{x_t^* - x_{t-1}^*} &= \norm{A_t^{-1} b_t - A_{t-1}^{-1} b_{t-1}} \\
            &= \norm{A_t^{-1} b_t - A_{t-1}^{-1} b_t
                                    + A_{t-1}^{-1} b_t - A_{t-1}^{-1} b_{t-1}} \\
            &\leq \norm{(A_t^{-1} - A_{t-1}^{-1}) b_t} + \norm{A_{t-1}^{-1} \xi_t} \\
            &\leq \norm{(A_t^{-1} - A_{t-1}^{-1})} \norm{b_t}
                + \norm{A_{t-1}} \xi_t
    \end{align*}

    We can bound this further by considering the suprenums of time-varying
    parameters. Define the following variables:
    \begin{align*}
        \Delta \bar{A} :&= \sup_{\tau > 1} \norm{A_t^{-1} - A_{t-1}^{-1}} \\
            &= \sup_{\tau > 1} \max_{0\leq i \leq n}
                \{|1 / a_{i,t} - 1 / a_{i,t-1}\} \\
        \bar{a} :&= \sup_{\tau > 0} {A_t^{-1}} \\
            &= \sup_{\tau > 0} \max_{0\leq i \leq n} \{
                |1 / a_{i,t}|
            \} \\
        \bar{b} :&= \sup_{\tau > 0} b_t \\
        \bar{\xi} :&= \sup_{\tau > 0} \xi_t
    \end{align*}
    
    Our final bound is then
    \begin{equation*}
        \norm{x_t^* - x_{t-1}^*} \leq \bar{a} \bar{\xi} + \Delta \bar{A} \bar{b}.
    \end{equation*}
\end{letterenum}

\section{Problem 2}

\begin{letterenum}
    \item We start by defining the maps $\hat{F}: \R^n \to \R^n$ and
    $F: \R^n \to \R^n$ such that:
    \begin{align*}
        \hat{F}_t(x) = \text{proj}_{\X} \{x - \eta \hat{\nabla}f_t(x)\} \\
        F_t(x) = \text{proj}_{\X} \{x - \eta \nabla f_t(x)\}
    \end{align*}
    Therefore the algorithmic map can be written $x_t = \hat{F}_t(x_{t-1})$.
    The optimum depends on the true gradient, rather than 
    the approximate, and is a fixed point of $F_t$. To see this, note that
    if $\text{proj}_\X \{x_t^* - \eta \nabla f_t(x_t^*)\} = x_t^*$, we have
    for any $y\in\X$
    \begin{align*}
        (x_t^* - x_t^* - \eta \nabla f_t(x_t^*))^T (y - x_t^*) \leq 0 \\
        \nabla f_t(x_t^*)^T (y - x_t^*) \geq 0 \\
    \end{align*}
    by [Theorem 9.8 Beck14]. $x_t^*$ is therefore a stationary point and
    the optimal solution of the problem $\min_{x\in\X} f_t(x)$ [Theorem 9.7 Beck].
    Therefore, $x_t^* = F_t(x_t^*)$.

    The single-step tracking bound is derived as follows:
    \begin{align*}
        \norm{x_t - x_t^*} &= \norm{\hat{F}_t(x_{t-1}) - x_t^*} \\
            &= \norm{\hat{F}_t(x_{t-1}) - F_t(x_{t-1}^*)
                                        + F_t(x_{t-1}^*) - F_t(x_t^*)} \\
            &\leq \norm{\hat{F}_t(x_{t-1}) - F_t(x_{t-1}^*)} +
                \norm{F_t(x_{t-1}^*) - F_t({x_t^*})} \\
            &\leq \rho \norm{x_{t-1} - x_{t-1}^*} + \eta \norm{e_t} + \rho
                \norm{x_{t-1}* - x_t^*} \\
            &= \rho \norm{x_{t-1} - x_{t-1}^*} + \eta \norm{e_t} + \rho \sigma_t \\
    \end{align*}

    where we have defined $\sigma_t = \norm{x_t - x_{t-1}}$. $\rho$ is the
    Lipschitz constant of the algorithmic map, again given by:
    \begin{equation*}
        \rho = \max\{|1 - \alpha \mu|, |1 - \alpha L|\}
    \end{equation*}

    The second inequality in the single step bound uses the following relation:
    \begin{align*}
        \norm{\hat{F}_t(x_{t-1}) - F_t(x_{t-1}^*)} &=
            \norm{\text{proj}_\X \{x_{t-1} - \eta \nabla f_t(x_{t-1}) - \eta e_t\} -
                \text{proj}_\X \{x_{t-1}^* - \eta \nabla f_t(x_{t-1}^*)\}} \\
            &\leq \norm{(x_{t-1} - \eta \nabla f_t(x) - \eta e_t) -
                (x_{t-1}^* - \eta \nabla f_t(x_{t-1}^*))} \\
            &\leq \norm{x_{t-1} - \eta \nabla f_t(x) - (x_{t-1}^* -
                \eta \nabla f_t(x_{t-1}^*))} + \eta \norm{e_t} \\
            &\leq \rho \norm{x_{t-1} - x_{t-1}^*} + \eta \norm{e_t}.
    \end{align*} 
    We have used the fact that
    $\hat{\nabla} f_t(x_{t-1}) = \nabla f_t(x_{t-1}) + e_t$ in the first
    equality, as well the nonexpansiveness of the
    projection operator in the first inequality.

    The multi-step bound proceeds by iterating upon the single-step bound:
    \begin{align*}
        \norm{x_t - x_t^*} &\leq \rho \biggl(
            \rho \norm{x_{t-2} - x_{t-2}^*} + \eta \norm{e_{t-1}} +
                \rho \sigma_{t-1}
        \biggr) \\
            &\qquad + \eta \norm{e_t} + \rho \sigma_t \\
        &= \rho^2 \norm{x_{t-2} - x_{t-2}^*} + \rho \eta \norm{e_{t-1}}
            + \rho^2 \sigma_{t-1} \\
            &\qquad + \eta \norm{e_t} + \rho \sigma_t \\
        &\leq \rho^t \norm{x_0 - x_0^*} + \sum_{i=1}^{t} \rho^{t-i} \eta \norm{e_i}
            + \sum_{i=1}^{t} \rho^{t-i+1} \sigma_i
    \end{align*}
    From this result, we see that the requested terms $D_t$ and $E_t$ are
    \begin{align*}
        D_t &= \sum_{i=1}^t \rho^{t-i+1} \sigma_i \\
        E_t &= \sum_{i=1}^t \rho^{t-i} \eta \norm{e_i} \\
    \end{align*}

    \item The suprenum in the limit as $t$ approaches infinity is the sum
    of the $\limsup$ of individual terms. Since each $\sigma_i < \bar{\sigma}$,
    we can pull $\bar{\sigma}$ outside of the summation.
    For $D_t$, we have
    \begin{align*}
        \limsup_{t\to\infty} D_t &\leq \limsup_{t\to\infty} \biggl(
            \sum_{i=1}^t \rho^{t-i+1} 
        \biggr) \bar{\sigma} \\
        &= \limsup_{t\to\infty} \biggl(
            \sum_{j=0}^t \rho^{j+1} 
        \biggr) \bar{\sigma} \\
        &\leq \frac{\rho \bar{\sigma}}{1 - \rho}.
    \end{align*}

    Similarly, for $E_t$, we have
    \begin{align*}
        \limsup_{t\to\infty} E_t &\leq \limsup_{t\to\infty} \biggl(.
            \sum_{i=1}^t \rho^{t-i} 
        \biggr) \eta \bar{e} \\
        &= \limsup_{t\to\infty} \biggl(
            \sum_{j=0}^t \rho^{j} 
        \biggr) \bar{\sigma} \\
        &\leq \frac{\bar{e}}{1 - \rho}.
    \end{align*}

    The $\rho^t \norm{x_0 - x_0^*}$ term goes to zero as $t \to\infty$
    since $\rho < 1$. Thus, the total $\limsup$ of the error has bound
    \begin{equation*}
        \limsup_{t\to\infty} \norm{x_t - x_t^*} \leq
            \frac{\rho \bar{\sigma}}{1 - \rho} + 
            \frac{ \bar{e}}{1 - \rho}.
    \end{equation*}



\end{letterenum}


\section{Problem 3}
\begin{letterenum}
    \item We check the Hessian of the function to verify convexity:
    \begin{align*}
        \nabla f([x_1, x_2]^T) &= \mat{
            \frac{1}{x_2} \\
            - \frac{x_1}{x_2^2}
        } \\
        \nabla^2 f([x_1, x_2]^T) &= \mat{
            0 & - \frac{1}{x_2} \\
            - \frac{1}{x_2} & \frac{2x_1}{x_2^3}
        } = \frac{1}{x_2^3} \mat{
            0 & - x_2^2 \\
            -x_2^2 & 2x_1 \\
        }
    \end{align*}

    The determinant of the hessian is $1/x_2^3 (0 - x_2^4) < 0$, since
    $x_2 > 0$ implies $x_2^3 > 0$ and $x_2^4 > 0$. Thus, the
    Hessian is not positive semidefinite [Proposition 2.20, Beck14],
    and therefore $f$ is not convex [Theorem 7.12, Beck14].

    \item Since $f$ is the piecewise maximum of two functions, it is
        convex if its interior functions are.

        As in Problem 1, the interior functions are convex if
        $A_1^T A_1 \succeq 0$ and $A_2^T A_2 \succeq 0$. Thus, $f$
        is convex when both these condition holds. If on the other
        hand $A_1^T A_1$ and $A_2^T A_2$ are both negative definite, $f$
        is not convex.

    \item We will show $f$ is not convex by counterexample. Take $p = 1/2$,
        $n=2$. Then, we have
        \begin{equation*}
            f(x) = (\sqrt{x_0} + \sqrt{x_1})^2
        \end{equation*}
        We will evaluate at $x = [2, 0]^T$ and $y=[0,2]^T$:
        \begin{align*}
            f(x) &= (\sqrt{2} + 0)^2 = 2 \\
            f(y) &= (0 + \sqrt{2})^2 = 2
        \end{align*}
        However, $f$ evaluated at the point $\lambda x + (1 - \lambda y)\in\R^2$,
        taking $\lambda=0.5$, is given by
        \begin{equation*}
            f(1/2 [2, 0]^T + 1/2 [0, 2]^T) = f([1, 1]^T)
                = (\sqrt{1} + \sqrt{1})^2 = 4.
        \end{equation*}
        while $\lambda f(x) + (1 - \lambda) f(y) = 2$.

        Thus, $f(\lambda x + (1 - \lambda)y) > \lambda f(x) + (1 - \lambda) f(y)$,
        and $f$ is not convex by the definition of a convex function
        [Def 7.1 Beck14].
\end{letterenum}


\end{document}
